% Question 3 formulas for paper handoff

\section{问题三：多工序生产的期望递推模型与遗传算法}

\subsection{生产结构与决策变量}

问题三包含 8 个零配件、3 个半成品和 1 个成品。生产结构为
\[
\{1,2,3\}\rightarrow H_1,\qquad
\{4,5,6\}\rightarrow H_2,\qquad
\{7,8\}\rightarrow H_3,
\]
\[
\{H_1,H_2,H_3\}\rightarrow F.
\]

定义 16 位二元决策向量
\[
X=(x_1,\ldots,x_{16}),\qquad x_i\in\{0,1\}.
\]

其中
\[
x_1,\ldots,x_8
\]
表示 8 个零配件是否检测；
\[
x_9,x_{10},x_{11}
\]
表示 3 个半成品是否检测；
\[
x_{12},x_{13},x_{14}
\]
表示检测出不合格半成品后是否拆解；
\[
x_{15}
\]
表示成品是否检测；
\[
x_{16}
\]
表示不合格成品是否拆解。

搜索空间大小为
\[
2^{16}=65536.
\]

\subsection{质量传递规则}

零配件没有主动报废或拆解决策。对零配件 \(i\)，若 \(x_i=1\)，则检测并丢弃检测出的不合格件，直到合格件进入装配；若 \(x_i=0\)，则不检测并直接进入装配。

拆解决策只针对不合格半成品或不合格成品。拆解不会损坏内部下级零部件，拆解后保留原来的真实质量状态，不重新按照次品率生成质量。

若任一输入不合格，则上级产品必为不合格。只有所有下级输入均合格时，上级产品才以对应装配次品率 \(p\) 额外产生装配缺陷。

\subsection{状态定义}

对每个零配件和半成品设置状态
\[
s\in\{N,KG,UG,UB\}.
\]
其中 \(N\) 表示当前没有该对象，\(KG\) 表示已知合格，\(UG\) 表示未检测但实际合格，\(UB\) 表示未检测且实际不合格。

系统状态记为
\[
S=(s_1,\ldots,s_8,h_1,h_2,h_3).
\]
初始状态为
\[
S_0=(N,\ldots,N).
\]

\subsection{零配件购买与检测}

若零配件 \(i\) 需要检测，则检测出坏件后丢弃，并重复购买和检测直到得到合格件。设零配件次品率为 \(p_i\)，购买单价为 \(c_i\)，检测成本为 \(t_i\)，则期望成本为
\[
C_i^{T}
=
\sum_{k=1}^{\infty}k(c_i+t_i)p_i^{k-1}(1-p_i)
=
\frac{c_i+t_i}{1-p_i}.
\]

若零配件 \(i\) 不检测，则购买成本为
\[
C_i^{N}=c_i,
\]
并以概率 \(1-p_i\) 进入状态 \(UG\)，以概率 \(p_i\) 进入状态 \(UB\)。

\subsection{半成品装配与检测}

设半成品 \(H_j\) 的下级零配件集合为 \(G_j\)，装配成本为 \(a_j\)，装配次品率为 \(p_{H_j}\)，检测成本为 \(T_j\)，拆解成本为 \(D_j\)。

若所有下级零配件实际合格，则
\[
P(H_j\text{ 合格})=1-p_{H_j},
\qquad
P(H_j\text{ 不合格})=p_{H_j}.
\]
若存在任一下级零配件不合格，则
\[
P(H_j\text{ 不合格})=1.
\]

若检测半成品且发现合格，则半成品进入状态 \(KG\)。若检测发现不合格，且选择拆解，则支付拆解成本 \(D_j\)，半成品消失，下级零配件按原真实质量状态回流；若不拆解，则该半成品及其下级零配件均被报废。

\subsection{成品装配、退换与拆解}

设成品装配成本为 \(a_F\)，装配次品率为 \(p_F\)，检测成本为 \(T_F\)，拆解成本为 \(D_F\)，市场售价为 \(R\)，调换损失为 \(L\)。

若三个半成品均实际合格，则
\[
P(F\text{ 合格})=1-p_F,
\qquad
P(F\text{ 不合格})=p_F.
\]
若存在任一半成品不合格，则
\[
P(F\text{ 不合格})=1.
\]

若不检测成品且成品不合格，则该产品进入市场并产生调换损失 \(L\)，退回后按 \(x_{16}\) 决定拆解或报废。若检测成品，则不合格品不会流入市场，不产生调换损失。

\subsection{固定策略的期望递推}

令
\[
V_X(S)
\]
表示在固定策略 \(X\) 下，从状态 \(S\) 出发最终交付 1 件合格成品所需的期望成本。对每个状态建立递推方程
\[
V_X(S)
=
C_X(S)
+
\sum_{S'}P_X(S'\mid S)V_X(S').
\]

将所有可达状态联立，得到
\[
(I-P_X)V_X=C_X.
\]

若矩阵奇异或无有限解，则说明该固定策略可能导致坏件拆解回流后反复进入装配，标记为不可行。否则初始状态期望成本为
\[
E[C\mid X]=V_X(S_0).
\]

期望利润为
\[
E[\Pi(X)]
=
R-V_X(S_0).
\]

问题三优化目标为
\[
X^\ast=\arg\max_X E[\Pi(X)].
\]

\subsection{遗传算法}

遗传算法染色体即为 16 位策略编码
\[
X=(x_1,\ldots,x_{16}).
\]

适应度函数取固定策略的解析期望利润
\[
\mathrm{Fitness}(X)=E[\Pi(X)].
\]

算法流程为随机初始化种群、锦标赛选择、单点交叉、按位变异、精英保留。程序采用参数
\[
N_{\mathrm{pop}}=100,\quad
G=200,\quad
p_c=0.8,\quad
p_m=0.02,\quad
N_{\mathrm{elite}}=5.
\]

为验证 GA 是否找到全局最优，程序同时枚举全部
\[
2^{16}=65536
\]
个策略，并与 GA 最优结果比较。
